\section{The inventory encoding and its reference surfaces}
\label{enc:section}

The inventory word determines the primal map together with an auxiliary
spanning tree and its complementary dual tree.  Their contours will
provide the reference sphere used in the topological comparison.
We first fix the time and height conventions.  One inventory letter represents
one exploration triangle, which consists of two of the flag triangles
of \Cref{intro:model-section}.  A word of length $2n$ has duration one
after division of time by $2n$.  Multiplying both contour coordinates by one
fixed constant makes their limiting marginal variances one.  This constant is
also used when contour heights are interpreted as quantum boundary lengths.

\subsection{The exact inventory law}

Write $\mathsf h,\mathsf c$ for the two burger types,
$\mathsf H,\mathsf C$ for their rigid orders, and $\mathsf F$ for a flexible
order.  Reading from left to right, an order removes the most recently
produced available burger of an allowed type.  Reduction retains the
unfulfilled orders, in chronological order, followed by the surviving stack.
Set
\begin{equation}\label{enc:letter-law}
 b=\frac{\sqrt q}{2+\sqrt q}\in(0,1/2),\qquad
 \PP(\mathsf h)=\PP(\mathsf c)=\frac14,\quad
 \PP(\mathsf H)=\PP(\mathsf C)=\frac{1-b}{4},\quad
 \PP(\mathsf F)=\frac b2.
\end{equation}
In the bilateral iid word every letter has a unique match almost surely
\cite[Proposition 2.2]{SheffieldInventory}.  Resolve $\mathsf F$ as
$\mathsf H$ or $\mathsf C$ according to its matched burger.  The resolved
contour $Z=(L,R)$ increases by the appropriate coordinate vector at a burger
and decreases by that vector at its order.  It is linearly interpolated
between integer times and anchored at $Z(0)=0$.  Write $N_{\mathsf F}$
for the number of flexible orders in a finite word.

\begin{proposition}[Inventory realization]\label{enc:inventory-law}
The iid word of length $2n$, conditioned to reduce to the empty word,
encodes exactly the rooted FK-decorated map law of the theorem.  Its
underlying primal map, with every primal edge retained, is the map obtained
by the tree/cotree sewing of its resolved contour.  The FK interface count
equals $1+N_{\mathsf F}$.
\end{proposition}

\begin{proof}
The inventory bijection is the rooted map bijection of
\cite[Sections 4.1--4.2]{SheffieldInventory}.  The conversion from a
spanning subgraph to the auxiliary spanning tree is the construction of
\citet[Section~4]{Bernardi2008Subgraphs}, as explained in
\citet[Section~4.1]{SheffieldInventory}.  Its elementary sewing objects
are incidence occurrences, so loops, bridges, and repeated incidences are
included.  A matched burger/order pair forms the two exploration triangles
in one quadrangle.  The $L$-pairs give the primal spanning-tree edges;
the $R$-pairs give the dual spanning-tree edges, whose quadrangles carry the
remaining primal edges.  Resolving a flexible order selects the auxiliary
tree/cotree diagonal.  Restoring its FK diagonal changes the decoration,
not the quadrangle or the underlying all-edge primal map.

A fully matched word has $n$ burgers and $n$ orders.  Its unconditioned
probability is
\[
 \left(\frac{1-b}{16}\right)^n
 \left(\frac{2b}{1-b}\right)^{N_{\mathsf F}}.
\]
The bijection's loop-opening operation gives one initial interface and one
additional interface for each flexible order.  Since
$2b/(1-b)=\sqrt q$, the probability is a constant depending only on $n,q$
times $q^{\#\Gamma/2}$, which is the specified finite map weight.
The rooted bijection is bijective, rather than a quotient weighted by
automorphism counts.  Uniform rerooting at an incidence preserves the law
because every map under consideration has the same number of root
incidences.  These observations also show why deleting cotree-associated
primal edges would change the metric in the theorem.
\end{proof}

\Cref{enc:flag-quadrangle} distinguishes the two exploration triangles
from the four flags of an original edge.  It also records the local
interface pairings determined by the edge's occupation status.

\begin{figure}[ht]
\centering
\begingroup
\begin{tikzpicture}[x=.92cm,y=.92cm,font=\small,
  boundary/.style={draw=black!65,line width=.6pt},
  seam/.style={draw=black!40,densely dashed,line width=.45pt},
  diagonal/.style={draw=linkblue,line width=1.2pt},
  interface/.style={draw=orange!85!black,line width=1.2pt}]

\begin{scope}[xshift=0cm]
  \fill[blue!9] (-1.45,0)--(0,1.12)--(0,0)--cycle;
  \fill[orange!12] (0,1.12)--(1.45,0)--(0,0)--cycle;
  \fill[blue!9] (1.45,0)--(0,-1.12)--(0,0)--cycle;
  \fill[orange!12] (0,-1.12)--(-1.45,0)--(0,0)--cycle;
  \draw[boundary] (-1.45,0)--(0,1.12)--(1.45,0)--(0,-1.12)--cycle;
  \draw[seam] (-1.45,0)--(1.45,0) (0,1.12)--(0,-1.12);
  \node[left] at (-1.45,0) {$v_0$};
  \node[right] at (1.45,0) {$v_1$};
  \node[above] at (0,1.12) {$f_0$};
  \node[below] at (0,-1.12) {$f_1$};
  \fill (0,0) circle (1.2pt);
  \node[fill=white,inner sep=1pt,above right] at (0,0) {$e$};
  \node at (0,1.90) {(a) Four flags};
\end{scope}

\begin{scope}[xshift=4.6cm]
  \fill[blue!9] (-1.45,0)--(0,1.12)--(1.45,0)--cycle;
  \fill[orange!12] (-1.45,0)--(0,-1.12)--(1.45,0)--cycle;
  \draw[boundary] (-1.45,0)--(0,1.12)--(1.45,0)--(0,-1.12)--cycle;
  \draw[seam] (0,1.12)--(0,-1.12);
  \draw[diagonal] (-1.45,0)--(1.45,0);
  \node[left] at (-1.45,0) {$v_0$};
  \node[right] at (1.45,0) {$v_1$};
  \node[above] at (0,1.12) {$f_0$};
  \node[below] at (0,-1.12) {$f_1$};
  \node at (0,1.90) {(b) Primal diagonal};
  \node[font=\scriptsize] at (0,-1.90) {two exploration triangles};
\end{scope}

\begin{scope}[xshift=9.2cm]
  \fill[blue!9] (-1.45,0)--(0,1.12)--(0,-1.12)--cycle;
  \fill[orange!12] (1.45,0)--(0,1.12)--(0,-1.12)--cycle;
  \draw[boundary] (-1.45,0)--(0,1.12)--(1.45,0)--(0,-1.12)--cycle;
  \draw[seam] (-1.45,0)--(1.45,0);
  \draw[diagonal] (0,1.12)--(0,-1.12);
  \node[left] at (-1.45,0) {$v_0$};
  \node[right] at (1.45,0) {$v_1$};
  \node[above] at (0,1.12) {$f_0$};
  \node[below] at (0,-1.12) {$f_1$};
  \node at (0,1.90) {(c) Dual diagonal};
  \node[font=\scriptsize] at (0,-1.90) {two exploration triangles};
\end{scope}

\begin{scope}[xshift=2.3cm,yshift=-4.4cm]
  \draw[boundary] (-1.45,0)--(0,1.12)--(1.45,0)--(0,-1.12)--cycle;
  \draw[interface] (-.97,.373) .. controls (-.55,.15) and (-.55,-.15) .. (-.97,-.373);
  \draw[interface] (.97,.373) .. controls (.55,.15) and (.55,-.15) .. (.97,-.373);
  \foreach \x/\y/\lab/\pos in {-.97/.373/1/above left,.97/.373/2/above right,.97/-.373/3/below right,-.97/-.373/4/below left}{
    \fill (\x,\y) circle (1.5pt);
    \node[\pos,inner sep=2pt,font=\scriptsize] at (\x,\y) {$p_{\lab}$};
  }
  \node at (0,1.75) {(d) Unoccupied edge};
\end{scope}

\begin{scope}[xshift=6.9cm,yshift=-4.4cm]
  \draw[boundary] (-1.45,0)--(0,1.12)--(1.45,0)--(0,-1.12)--cycle;
  \draw[interface] (-.97,.373) .. controls (-.3,.62) and (.3,.62) .. (.97,.373);
  \draw[interface] (-.97,-.373) .. controls (-.3,-.62) and (.3,-.62) .. (.97,-.373);
  \foreach \x/\y/\lab/\pos in {-.97/.373/1/above left,.97/.373/2/above right,.97/-.373/3/below right,-.97/-.373/4/below left}{
    \fill (\x,\y) circle (1.5pt);
    \node[\pos,inner sep=2pt,font=\scriptsize] at (\x,\y) {$p_{\lab}$};
  }
  \node at (0,1.75) {(e) Occupied edge};
\end{scope}
\end{tikzpicture}
\endgroup

\caption{Incidences at one original edge.  The four flags have labels
$(v_i,e,f_j)$, $i,j\in\{0,1\}$.  Selecting the primal or dual auxiliary
diagonal groups them into the two exploration triangles of the matched
letters.  The lower panels show the two pairings of the four interface
ports: an unoccupied edge pairs ports about each primal endpoint,
whereas an occupied edge pairs them about each incident face.
This is a topological schematic, not an isometric drawing of the
equilateral flags.  Labels may repeat for a loop or bridge, but the
four flag and port occurrences remain distinct.}
\label{enc:flag-quadrangle}
\end{figure}

\begin{lemma}[Tree and incidence identities]\label{enc:incidence}
For a fully matched resolved word, let $T_n$ be its primal spanning tree
and $T_n^*$ the complementary dual spanning tree.  At the cut after
letter $k$, the exploration has a primal corner vertex $v_n(k)$ and a
dual corner face $f_n(k)$.  In unscaled integer height units,
\begin{align*}
 d_{T_n}(v_n(k),v_n(l))
 &=L_k+L_l-2\min_{j\in[k,l]}L_j,\\
 d_{T_n^*}(f_n(k),f_n(l))
 &=R_k+R_l-2\min_{j\in[k,l]}R_j,
 \qquad 0\le k\le l\le2n.
\end{align*}
Every primal vertex and dual face occurrence is represented by these
corner sequences.  Each quadrangle is the union of its two matched
exploration triangles; each such triangle is incident to the corner
states immediately before and after its letter.  Thus the endpoints,
faces, and flag incidences of a quadrangle have a uniformly bounded
collection of tour representatives.  The representatives for each
primal vertex are identified by the $L$-tree rule, and those for each
dual vertex by the $R$-tree rule; moving from a cut to its adjacent
triangle changes a height by at most one.
\end{lemma}

\begin{proof}
After resolution this is the ordinary tree/cotree contour exploration
in \cite[Section 4.1, Figures 6--7]{SheffieldInventory}.  A burger
creates an edge to a child in its tree and its order returns along
that edge.  Letters of the other type hold this tree position fixed.
The depth of the common ancestor of the positions at $k,l$ is the
minimum contour height between them, proving the two distance
identities.  The exploration crosses each tree-edge side and each
quadrangle side occurrence, which proves coverage of the corners,
including repeated occurrences.  A letter triangle is bounded by the
two successive exploration sides and the created or traversed tree
edge.  Pairing it with its matching letter triangle gives exactly the
associated quadrangle.  Subdivision into flags introduces an edge
midpoint but no additional incidence choice.  This proves the last
assertions without any bound on vertex or face degrees.
\end{proof}

\begin{theorem}[Contour limits]\label{enc:contour-limits}
Let
\[
 \varrho=\frac b{1-b}=-\cos\left(\frac{\pi\gamma^2}{4}\right),
 \qquad \gamma\in(\sqrt2,2),\qquad \kappa=\frac{16}{\gamma^2}.
\]
The diffusively rescaled bilateral resolved contours converge, uniformly
on compact time intervals in law, to a bilateral Brownian motion whose
marginal variance rates are one and whose correlation is $\varrho$.
For a length-$2n$ word conditioned to reduce to the empty word, the same
rescaling, with duration one, converges uniformly in law to the corresponding
quadrant Brownian excursion from $(0,0)$ to $(0,0)$.

For an unconditioned raw block of $O(n)$ letters, the number of flexible
orders left unidentified by reduction of the block is $o_{\PP}(\sqrt n)$.
Consequently changing their resolutions changes the rescaled contour
increments by $o_{\PP}(1)$ uniformly on that block.
\end{theorem}

\begin{proof}
The bilateral assertion is \cite[Theorem 2.5]{SheffieldInventory}.
In the original units the covariance rates of the two net burger counts
are $(1-b)/2,(1-b)/2,b/2$; multiplying heights by
$\sqrt{2/(1-b)}$ gives the stated convention.  The conditioned assertion is
\cite[Theorem 1.8]{GwynneSunFKIII}, followed by the deterministic change
from duration two to duration one.  The unidentified-order estimate is
\cite[Remark 3.17]{SheffieldInventory}; see also the reduced-word estimates
in \cite{GwynneMaoSunFKI}.  Each changed resolution changes a contour
coordinate by at most one, so their total number bounds the supremum error.
This last assertion concerns a raw unconditioned block.  The transfer of
discrete observations to conditioned extractions will be proved in
Section~\ref{ker:section}.
\end{proof}

\subsection{The reference LQG surfaces}

\begin{theorem}[Peanosphere realization]\label{enc:peanosphere}
The bilateral Brownian contour is the boundary-length process of a
$\gamma$-quantum cone explored by an independent, unparameterized
whole-plane space-filling $\mathrm{SLE}_{\kappa}$, subsequently
parameterized by quantum area.  The duration-one excursion is the
boundary-length process of the ordinary unit-area $\gamma$-quantum sphere
explored in the same way from an area-typical marked point $p_*$.
The contour determines the curve-decorated quantum surface up to its
marked conformal equivalence.

Write $\eta$ for the exploration parameterized by quantum area.
In either realization, $\eta_*(\dd t)=\mu_h$.  Its spatial identifications
are generated by finite chains of scalar horizontal chords
\begin{equation}\label{enc:chord-rule}
 s\sim_L t\quad\Longleftrightarrow\quad
 L_s=L_t=\inf_{[s,t]}L,
 \qquad s<t,
\end{equation}
and the analogous $R$-chords, with the usual endpoint identification for
the sphere.  The resulting quotient carries the topology of the plane,
or sphere, respectively.  Frontier arcs are parameterized by quantum
length, with running-minimum contour heights as their length coordinates.
\end{theorem}

\begin{proof}
For the cone, the surface and contour identification are the
mating-of-trees and welding theorems
\cite[Theorems~1.9 and~1.11]{DMSMating}. For the ordinary area-one sphere,
\cite[Theorem~1.1]{MillerSheffieldSphere} supplies the Brownian-loop law
and determination of the surface explored by an independent unclocked
whole-plane space-filling SLE. In our range $\gamma^2\in(2,4)$ the
covariance is the one in these theorems. The marked-point sampling rule
is \cite[Proposition~A.13]{DMSMating}. We use the common deterministic
boundary-length convention specified above, with area-clock duration one.

Call a scalar chord \emph{strict} if the contour is strictly above its
endpoint height between the endpoints.  The finite-chain assertion uses
a further geometric part of the DMS
construction. There are finitely many visits to each spatial point;
ordering their left and right flow-line sectors links consecutive visits
by scalar chords. This is the smallest generated equivalence relation in
\cite[Lemmas~8.9, 8.13--8.15 and Theorem~8.18]{DMSMating}, without taking
its closure. A non-strict scalar chord can be split at its possible
intermediate minimum visit into adjacent strict chords; scalar fibers
have at most three visits. The same finite-visit and adjacent-sector
argument applies to the independent whole-plane exploration of the
sphere away from its exploration root: changing to quantum-area time
preserves the visits, their order and their frontier incidences. At the
root one additionally identifies $0$ and $1$. Thus the sphere statement
combines its Brownian-loop realization with this geometric argument;
it does not invoke a separately stated sphere finite-chain theorem.
The quotient is topological because the resulting continuous bijection
from the compact time quotient to the sphere is a homeomorphism.
Running-minimum frontier coordinates are the quantum lengths supplied
by the same welding construction.
\end{proof}

We use actual field heights when comparing passage laws.  The following
description records explicitly the area normalization that couples
different patches.

\begin{proposition}[Three-mark sphere representation]
\label{enc:sphere-representation}
Let $h_{\C}$ be the whole-plane Gaussian free field (GFF) with circle average zero on
$\partial\mathbb D$, and put
\[
 |z|_+=\max\{1,|z|\},\qquad
 G_{\C}(z,w)=\log\frac{|z|_+|w|_+}{|z-w|},\qquad
 Q=\frac2\gamma+\frac\gamma2.
\]
Define
\begin{equation}\label{enc:liouville-field}
 h_L=h_{\C}+\gamma G_{\C}(0,\cdot)+\gamma G_{\C}(1,\cdot)
       -(2Q-\gamma)\log|\cdot|_+,
 \quad A=\mu_{h_L}(\C).
\end{equation}
Tilt the law of $h_L$ by the integrable strictly positive density
$A^{4/\gamma^2-2}/\EE[A^{4/\gamma^2-2}]$, and then set
\begin{equation}\label{enc:area-normalization}
 h=h_L-\gamma^{-1}\log A.
\end{equation}
The result is the ordinary unit-area sphere with three independent
area marks embedded at $\infty,0,1$.  Resampling its three area marks and
renormalizing gives the same law.  In particular, on any bounded region
avoiding $0,1$, its field modulo constants is locally mutually absolutely
continuous with a whole-plane GFF modulo constants.  Near $0$, away from
$1$, the corresponding comparison field is $h_{\C}-\gamma\log|\cdot|$.
\end{proposition}

\begin{proof}
This fixed-area field law originates in the sphere construction of
David, Kupiainen, Rhodes and Vargas
\citep[Section~3.3, Proposition~3.8 and Section~3.4]{DKRVSphere2016}.
We use its form in \citet[Definition~2.2]{BorgaGwynneSunPermutons};
\citet[Theorem~1.1 and Remark~1.2]{AruHuangSunSphere} identify it with
the ordinary DMS sphere.
The mark resampling statement and local comparison, including the
$\gamma$-insertion, are \cite[Lemmas 2.3--2.4]{BorgaGwynneSunPermutons}.
Subtracting the constant in \eqref{enc:area-normalization} gives area one
by the identity $\mu_{h+c}=e^{\gamma c}\mu_h$.
\end{proof}

The proposition deliberately retains the common random constant
$-\gamma^{-1}\log A$.  Local equivalence modulo constants alone does not
give a joint law for actual heights in several patches.  The required
joint density is derived from \eqref{enc:liouville-field}--
\eqref{enc:area-normalization} in Section~\ref{ker:section}.
